% ==================== THEME 3 — CONDITIONS ====================

% ---- cond-egalite ----
\element{bank-cond-egalite}{
    \begin{question}{cond-egalite-v1}
        Quel opérateur permet de \textbf{tester} l'égalité entre deux entiers (et non affecter) ?
        \begin{multicols}{4}
            \begin{reponses}
                \mauvaise{\texttt{=}}
                \bonne{\texttt{==}}
                \mauvaise{\texttt{=<}}
                \mauvaise{\texttt{:=}}
            \end{reponses}
        \end{multicols}
    \end{question}
}
\element{bank-cond-egalite}{
    \begin{question}{cond-egalite-v2}
        Quel opérateur permet de \textbf{tester} l'égalité entre deux caractères (et non affecter) ?
        \begin{multicols}{4}
            \begin{reponses}
                \mauvaise{\texttt{=}}
                \bonne{\texttt{==}}
                \mauvaise{\texttt{=<}}
                \mauvaise{\texttt{:=}}
            \end{reponses}
        \end{multicols}
    \end{question}
}
\element{bank-cond-egalite}{
    \begin{question}{cond-egalite-v3}
        Quel opérateur permet de \textbf{tester} l'égalité entre deux variables flottantes (et non affecter) ?
        \begin{multicols}{4}
            \begin{reponses}
                \mauvaise{\texttt{=}}
                \bonne{\texttt{==}}
                \mauvaise{\texttt{=<}}
                \mauvaise{\texttt{:=}}
            \end{reponses}
        \end{multicols}
    \end{question}
}
\element{conditions}{
    \restituegroupe[1]{bank-cond-egalite}
}

% ---- cond-intervalle ----
\element{bank-cond-intervalle}{
    \begin{question}{cond-intervalle-v1}
        Comment écrire la condition \(5 \le x \le 10\) (inclus) ?
        \begin{multicols}{4}
            \begin{reponses}
                \bonne{\lstinline[language=C]|x >= 5 && x <= 10|}
                \mauvaise{\lstinline[language=C]!x > 5 || x < 10!}
                \mauvaise{\lstinline[language=C]|x > 5 && x < 10|}
                \mauvaise{\lstinline[language=C]|x <= 5 && x >= 10|}
            \end{reponses}
        \end{multicols}
    \end{question}
}
\element{bank-cond-intervalle}{
    \begin{question}{cond-intervalle-v2}
        Comment écrire la condition \(0 \le n \le 20\) (inclus) ?
        \begin{multicols}{4}
            \begin{reponses}
                \bonne{\lstinline[language=C]|n >= 0 && n <= 20|}
                \mauvaise{\lstinline[language=C]!n > 0 || n < 20!}
                \mauvaise{\lstinline[language=C]|n > 0 && n < 20|}
                \mauvaise{\lstinline[language=C]|n <= 0 && n >= 20|}
            \end{reponses}
        \end{multicols}
    \end{question}
}
\element{bank-cond-intervalle}{
    \begin{question}{cond-intervalle-v3}
        Comment écrire la condition \(1 \le note \le 6\) (inclus) ?
        \begin{multicols}{4}
            \begin{reponses}
                \bonne{\lstinline[language=C]|note >= 1 && note <= 6|}
                \mauvaise{\lstinline[language=C]!note > 1 || note < 6!}
                \mauvaise{\lstinline[language=C]|note > 1 && note < 6|}
                \mauvaise{\lstinline[language=C]|note <= 1 && note >= 6|}
            \end{reponses}
        \end{multicols}
    \end{question}
}
% v4 : variante supplémentaire demandée explicitement pour l'intervalle
% 5 <= x <= 10 ; la variable est renommée (\texttt{score}) pour ne pas
% dupliquer mot pour mot la v1, qui utilise déjà les mêmes bornes sur \texttt{x}.
\element{bank-cond-intervalle}{
    \begin{question}{cond-intervalle-v4}
        Comment écrire la condition \(5 \le score \le 10\) (inclus) ?
        \begin{multicols}{4}
            \begin{reponses}
                \bonne{\lstinline[language=C]|score >= 5 && score <= 10|}
                \mauvaise{\lstinline[language=C]!score > 5 || score < 10!}
                \mauvaise{\lstinline[language=C]|score > 5 && score < 10|}
                \mauvaise{\lstinline[language=C]|score <= 5 && score >= 10|}
            \end{reponses}
        \end{multicols}
    \end{question}
}
\element{conditions}{
    \restituegroupe[1]{bank-cond-intervalle}
}

% ---- cond-si-erreur-frequente ----
\element{bank-cond-si-erreur-frequente}{
    \begin{question}{cond-si-erreur-frequente-v1}
        Quelle est l'erreur dans le code suivant ?
        \lstinputlisting{sources/codes/cond-erreur-egal.c}
        \begin{multicols}{2}
            \begin{reponses}
                \bonne{\texttt{=} affecte 5 à x au lieu de tester l'égalité (il fallait \texttt{==})}
                \mauvaise{Il manque un point-virgule après le \texttt{if}}
                \mauvaise{Les accolades sont interdites après un \texttt{if}}
                \mauvaise{\texttt{printf} ne peut pas être utilisé dans un \texttt{if}}
            \end{reponses}
        \end{multicols}
    \end{question}
}
\element{bank-cond-si-erreur-frequente}{
    \begin{question}{cond-si-erreur-frequente-v2}
        Quelle est l'erreur dans le code suivant ?
        \lstinputlisting{sources/codes/cond-erreur-egal-2.c}
        \begin{multicols}{2}
            \begin{reponses}
                \bonne{\texttt{=} affecte 10 à note au lieu de tester l'égalité (il fallait \texttt{==})}
                \mauvaise{Il manque un point-virgule après le \texttt{if}}
                \mauvaise{Les accolades sont interdites après un \texttt{if}}
                \mauvaise{\texttt{printf} ne peut pas être utilisé dans un \texttt{if}}
            \end{reponses}
        \end{multicols}
    \end{question}
}
\element{bank-cond-si-erreur-frequente}{
    \begin{question}{cond-si-erreur-frequente-v3}
        Quelle est l'erreur dans le code suivant ?
        \lstinputlisting{sources/codes/cond-erreur-egal-3.c}
        \begin{multicols}{2}
            \begin{reponses}
                \bonne{\texttt{=} affecte 1 à statut au lieu de tester l'égalité (il fallait \texttt{==})}
                \mauvaise{Il manque un point-virgule après le \texttt{if}}
                \mauvaise{Les accolades sont interdites après un \texttt{if}}
                \mauvaise{\texttt{printf} ne peut pas être utilisé dans un \texttt{if}}
            \end{reponses}
        \end{multicols}
    \end{question}
}
\element{conditions}{
    \restituegroupe[1]{bank-cond-si-erreur-frequente}
}

% ---- cond-else-if ----
\element{bank-cond-else-if}{
    \begin{question}{cond-else-if-v1}
        Dans une structure \texttt{if}/\texttt{else if}/\texttt{else}, combien de blocs au maximum peuvent s'exécuter pour un même passage ?
        \begin{multicols}{4}
            \begin{reponses}
                \bonne{1}
                \mauvaise{Autant que de conditions vraies}
                \mauvaise{2}
                \mauvaise{0}
            \end{reponses}
        \end{multicols}
    \end{question}
}
\element{bank-cond-else-if}{
    \begin{question}{cond-else-if-v2}
        Dans une structure \texttt{if}/\texttt{else if}/\texttt{else if}/\texttt{else}, combien de blocs au maximum sont exécutés pour un même passage ?
        \begin{multicols}{4}
            \begin{reponses}
                \bonne{1}
                \mauvaise{Autant que de conditions vraies}
                \mauvaise{2}
                \mauvaise{0}
            \end{reponses}
        \end{multicols}
    \end{question}
}
\element{bank-cond-else-if}{
    \begin{question}{cond-else-if-v3}
        Dans une structure \texttt{if}/\texttt{else}, combien de blocs au maximum peuvent s'exécuter pour un même passage ?
        \begin{multicols}{4}
            \begin{reponses}
                \bonne{1}
                \mauvaise{Autant que de conditions vraies}
                \mauvaise{2}
                \mauvaise{0}
            \end{reponses}
        \end{multicols}
    \end{question}
}
\element{conditions}{
    \restituegroupe[1]{bank-cond-else-if}
}

% ---- cond-et-ou ----
% Durci : un simple "&&" à deux termes se résout par calcul direct sans
% connaître les priorités d'opérateurs. On passe à un mélange && / || à trois
% termes où \&\& est prioritaire sur \textbar\textbar : une lecture naïve
% gauche-à-droite (comme si les deux opérateurs avaient la même priorité)
% donne un résultat différent du résultat correct.
\element{bank-cond-et-ou}{
    \begin{question}{cond-et-ou-v1}
        Que vaut l'expression \lstinline[language=C]!(3 > 1) || (2 > 5) && (4 > 10)! ?
        \begin{multicols}{5}
            \begin{reponses}
                \bonne{Vrai (1)}
                \mauvaise{Faux (0)}
                \mauvaise{Erreur de compilation}
                \mauvaise{0}
                \mauvaise{10}
            \end{reponses}
        \end{multicols}
    \end{question}
}
\element{bank-cond-et-ou}{
    \begin{question}{cond-et-ou-v2}
        Que vaut l'expression \lstinline[language=C]!(5 > 2) || (1 > 9) && (7 > 20)! ?
        \begin{multicols}{5}
            \begin{reponses}
                \bonne{Vrai (1)}
                \mauvaise{Faux (0)}
                \mauvaise{Erreur de compilation}
                \mauvaise{0}
                \mauvaise{20}
            \end{reponses}
        \end{multicols}
    \end{question}
}
\element{bank-cond-et-ou}{
    \begin{question}{cond-et-ou-v3}
        Que vaut l'expression \lstinline[language=C]!(6 > 4) || (0 > 8) && (3 > 15)! ?
        \begin{multicols}{5}
            \begin{reponses}
                \bonne{Vrai (1)}
                \mauvaise{Faux (0)}
                \mauvaise{Erreur de compilation}
                \mauvaise{0}
                \mauvaise{15}
            \end{reponses}
        \end{multicols}
    \end{question}
}
\element{conditions}{
    \restituegroupe[1]{bank-cond-et-ou}
}

% ---- cond-ecrire (écriture de code : reste \evaluationProfGrille, pas un tableau) ----
\element{bank-cond-ecrire}{
    \begin{question}{cond-ecrire-v1}
        Écrire une condition en langage C qui affiche \texttt{"Majeur"} si la variable entière \texttt{age} est supérieure ou égale à 18, et \texttt{"Mineur"} sinon.
        \evaluationProfGrille{2}{9}
    \end{question}
}
\element{bank-cond-ecrire}{
    \begin{question}{cond-ecrire-v2}
        Écrire une condition en langage C qui affiche \texttt{"Positif"} si la variable entière \texttt{n} est strictement supérieure à 0, et \texttt{"Négatif ou nul"} sinon.
        \evaluationProfGrille{2}{9}
    \end{question}
}
\element{bank-cond-ecrire}{
    \begin{question}{cond-ecrire-v3}
        Écrire une condition en langage C qui affiche \texttt{"Admis"} si la variable entière \texttt{note} est supérieure ou égale à 10, et \texttt{"Recalé"} sinon.
        \evaluationProfGrille{2}{9}
    \end{question}
}
\element{conditions}{
    \restituegroupe[1]{bank-cond-ecrire}
}
